\section{Agents as Data Researchers: Identifying Challenges in Data-Centric RSI}
\label{sec:evidencegaps}

We instantiate data-centric RSI as an iterative environment-generation and
reinforcement-learning loop. An agent serves as the data researcher,
constructing executable environments and tasks in which a target model trains
with online RL. Feedback from training and evaluation informs the next round.
We first formalize this process, then examine challenges in its research
context and generated training instances.

\textbf{Formalizing the data-research loop.}\quad
A \emph{research context} is the information available to an agent for
research decisions, encompassing the research history, current state, and
other decision-relevant information. Let $C_t$ denote this context at round
$t$ and $\operatorname{DataResearcher}$ the agent that uses it to construct
training instances. Each instance is a triple $(e,x,v)$, consisting of an
executable environment $e$, a task $x$ specifying a goal and initial
conditions, and a verifier $v$ assessing task completion. The agent's output
is a set $D_t$ of $n_t$ such instances:
% Keep these three displays readable even after short lead-in lines.
% This local spacing does not change the conference style or other equations.
\begingroup
\setlength{\abovedisplayskip}{8pt plus 2pt}
\setlength{\belowdisplayskip}{8pt plus 2pt}
\setlength{\abovedisplayshortskip}{8pt plus 2pt}
\setlength{\belowdisplayshortskip}{8pt plus 2pt}
\begin{equation}
\label{eq:researchoutput}
D_t=\{(e_j,x_j,v_j)\}_{j=1}^{n_t}
   =\operatorname{DataResearcher}(C_t).
\end{equation}
The set $D_t$ is then submitted for training and evaluation. Let $\pi_t$
denote the target model at the start of round $t$ and
$\operatorname{TrainEval}$ the fixed online-RL and public-evaluation
procedure. It returns the updated model $\pi_{t+1}$ and feedback $F_t$
comprising training records and evaluation results:
\begin{equation}
\label{eq:researchloop}
(\pi_{t+1},F_t)
= \operatorname{TrainEval}(\pi_t,D_t).
\end{equation}
Finally, $\operatorname{UpdateContext}$ integrates the submitted set and its
feedback with the prior research context to guide the next round:
\begin{equation}
\label{eq:contextupdate}
C_{t+1}=\operatorname{UpdateContext}(C_t,D_t,F_t).
\end{equation}
\endgroup

In this setting, the loop starts from an empty environment library. Guided
by feedback and the evolving research context, the data researcher repeatedly
creates new environments or revises existing ones, with the goal of improving
the target model over successive training rounds.

\subsection{Incomplete Behavioral Analysis}
\label{sec:coveragegap}

The research context must preserve task requirements and model behavior
to support training-target selection. To examine what survives this
analysis, we give an analyst agent (Astra/high) complete access to the
initial public-eval trajectories of Qwen3-8B reasoning: 1,200 BFCL, 417
$\tau^2$, and 75 ACEBench-Agent trajectories. The instruction requests an
analysis and specifies report fields, while leaving the reading strategy to
the agent. We measure report coverage of source-grounded task requirements and
failure mechanisms, crediting full preservation of their consequential
conditions and dependencies (Section~\ref{sec:analysiscoverage}).

The analysis reveals a substantial gap: task-requirement coverage
is 61.5--82.5\%, while failure-mechanism coverage ranges from only 39.4\% on
BFCL to 67.5\% on ACE. These automatic, task-averaged scores use the same final
paired evaluation as Section~\ref{sec:analysiscoverage}; the remainder includes
both partial descriptions and wholly missing mechanisms. Full-corpus access
thus leaves substantial behavioral information incompletely represented.
Coverage is lower on the larger corpora in this comparison, although benchmark
composition and corpus size vary together, preventing a causal size claim.

The research context therefore needs to preserve the mechanisms and conditions
that matter for the next data decision, with evidence available for inspection.

\subsection{Disconnected Data, Training, and Evaluation}
\label{sec:feedbackgap}

Updating the research context requires connecting intended training targets,
tasks actually trained, and subsequent evaluation outcomes. In ACE R9, one
revision contributed zero of its ten planned tasks. We further
audited 28 iterations of the original 4B studies, measuring the fraction of
trained tasks with nonconstant rewards across eight valid rollouts---the
within-task contrast used by GRPO. Among complete, valid groups, the
task-weighted rates were 74.6\%, 42.2\%, and 37.2\% on $\tau^2$, BFCL,
and ACE. These checks distinguish
whether an intervention reached training and whether its rollouts supplied
reward contrast, motivating explicit linkage to its intended target and
evaluation outcomes.
Section~\ref{sec:analysissignal} compares the per-round signal with
\texttt{Codex-Vanilla}.

\subsection{Incomplete Environment Implementation}
\label{sec:implementationgap}

A retail-exchange environment required matching identical items to different
replacements, yet all six inspected tasks assigned the same replacement to
both items. Checking the generator against the specification confirmed that
the required branch was absent.
This exposes a gap in constructing the output triples: a runnable instance
need not realize the intended training mechanism. Auditing environment
behavior, task requirements, and verifier scores tests whether the three
components jointly deliver the training target.
\todo{Cross-benchmark audit: 20 environments each on $\tau^2$, BFCL, and ACE;
eight rollouts per environment; LLM judgments of target implementation,
interaction, and verifier defects. Insert observed defect rates.}
